\clearpage
\section{Long Lagrangian}
\begin{colbox}
    The following Lagrangian is given 
    \begin{equation}
        \lagr = \lagr_{EM}+\lagr_k+\lagr_{int} = -\frac{1}{16\pi c}\tensor{F}{_\mu_\nu}\tensor{F}{^\mu^\nu}+\frac{k^2}{8\pi c}A_\mu A^\mu - \frac{1}{c^2}J_\mu A^\mu
    \end{equation}
    \begin{enumerate}
        \item What are the units of \(k\) and \(ck\)?
        \item Is the new action A lorentz scalar?
        \item Is it invariant for gauge transformations?
        \item Find the equations of motion
        \item Apply divergence to both sides and show the continuity equation does not hold unless \(\partial_\mu A^\mu=0\)
        \item Using the last part show that the equations of motion can be written as
        \begin{equation}
            \ins{\dalembertian+k^2}A^\mu = \frac{4\pi}{c}J^\mu
        \end{equation}
        \item Show that in the static limit for a point charge \(q\) at the origin the solution is
        \begin{equation}
            \phi(r) = \frac{q}{r}e^{-kr}
        \end{equation}
    \end{enumerate}
\end{colbox}
\subsection{Units}
Since all elements must have the same units we can demand that
\begin{equation}
    \left[-\frac{1}{16\pi c}\tensor{F}{_\mu_\nu}\tensor{F}{^\mu^\nu}\right] = \left[\frac{k^2}{8\pi c}A_\mu A^\mu\right]
\end{equation}
we know that 
\begin{equation}
    \left[\tensor{F}{_\mu_\nu}\right] = \left[\partial_\mu A_\nu\right] = \frac{\left[A\right]}{L}
\end{equation}
so 
\begin{equation}
    \frac{\left[A\right]^2}{L^2}=\left[k\right]^2\left[A\right]^2 \implies \left[k\right] = \frac{1}{L},\quad\left[ck\right]=\frac{1}{T}
\end{equation}
\subsection{Lorentz Scalar}
\(\tensor{F}{_\mu_\nu}\tensor{F}{^\mu^\nu}\) is known to be a lorentz scalar so we continue to the other elements.
\begin{align}
    A_\mu A^\mu \rightarrow \ins{\tensor{\Lambda}{_\mu^\nu}A_\nu}\ins{\tensor{\Lambda}{^\mu_\sigma}A^\sigma} = \tensor{\Lambda}{_\mu^\nu}\tensor{\Lambda}{^\mu_\sigma}A_\nu A^\sigma = \delta^\nu_\sigma A_\nu A^\sigma = A^\nu A_\nu
\end{align}
therefore this element is a lorentz scalar. As for the last element the process is identical to this entirely just with the name J instead of one of the As so I won't rewrite it. In total, the action is a lorentz scalar.
\subsection{Gauge Invariance}
A gauge transofmation is 
\begin{equation}
    A_\mu \rightarrow A_\mu+\partial_\mu\Lambda
\end{equation}
where \(\Lambda\) is a function of the coordinates. Looking at the kinetic element
\begin{align}
    A_\mu A^\mu &\rightarrow \ins{A_\mu+\partial_\mu\Lambda}\ins{A^\mu+\partial^\mu\Lambda}\\
    &= A_\mu A^\mu + A_\mu\partial^\mu\Lambda + A^\mu\partial_\mu\Lambda + \partial_\mu\Lambda \partial^\mu\Lambda
\end{align}
which does not automatically always equal \(A_\mu A^\mu\), nor is it a total derivative, so this is not gauge invariant, making the entire Lagrangian not gauge invariant.
\subsection{Equations of Motion}
Looking at a subset of the lagrangian, \(\lagr_{int}+\lagr_{EM}\) we have a known problem we solved in the lecture whose equation of motion are 
\begin{equation}
    \partial^\nu \tensor{F}{_\nu_\mu} = \frac{4\pi}{c}J_\mu
\end{equation}
therefore we only need to solve for the kinetic element
\begin{align}
    \pdv{\lagr_k}{A^\nu} &= \frac{k^2}{8\pi c}\ins{\pdv{A_\mu}{A^\nu}A^\mu + \pdv{A^\mu}{A^\nu}A_\mu}\\
    &= \frac{k^2}{8\pi c}\ins{\pdv{A^\sigma}{A^\nu}g_{\mu\sigma} A^\mu+A_\mu\delta^\mu_\nu}\\
    &= \frac{k^2}{8\pi c}\ins{\delta^\sigma_\nu g_{\mu\sigma} A^\mu+A_\mu\delta^\mu_\nu}\\
    &= \frac{k^2}{4\pi c}A_\nu
\end{align}
and 
\begin{equation}
    \pdv{\lagr_k}{\partial_nu A_\mu}=0
\end{equation}
so the equations of motion are 
\begin{equation}
    \partial^\nu \tensor{F}{_\nu_\mu} + k^2 A_\mu = \frac{4\pi}{c}J_\mu
\end{equation}
\subsection{Divergence}
directly
\begin{align}
    \partial^\mu\partial^\nu \tensor{F}{_\nu_\mu} + \partial^\mu k^2 A_\mu = \partial^\mu\frac{4\pi}{c}J_\mu
\end{align}
but since \(\partial^\mu\partial^\nu\) is symmetric and \(\tensor{F}{_\mu_\nu}\) is antisymmetric this element is zero.
\begin{align}
    \partial^\mu k^2 A_\mu = \partial^\mu\frac{4\pi}{c}J_\mu
\end{align}
and we can directly see that for the continuity equation to hold, we must have \(\partial^\mu A_\mu=0\).
\subsection{Box}
returning to the equation of motion and opening terms in it 
\begin{align}
    \partial^\nu \tensor{F}{_\nu_\mu} + k^2 A_\mu &= \frac{4\pi}{c}J_\mu\\
    \partial^\nu \ins{\partial_\nu A_\mu - \partial_\mu A_\nu} + k^2 A_\mu &= \frac{4\pi}{c}J_\mu\\
    \ins{\partial_\nu\partial^\nu A_\mu - \partial_\mu\partial^\nu A_\nu} + k^2 A_\mu &= \frac{4\pi}{c}J_\mu\\
    \ins{\dalembertian A_\mu - 0} + k^2 A_\mu &= \frac{4\pi}{c}J_\mu\\
    \ins{\dalembertian+k^2} A_\mu &= \frac{4\pi}{c}J_\mu
\end{align}
\subsection{Static limit}
Since we are in the static limit, the 4-current is of the form 
\begin{equation}
    j^\mu = \ins{c\rho,\vb{0}}
\end{equation}
where 
\begin{equation}
    \rho = q\delta^3(r)
\end{equation}
inserting into the last part and looking at the time component 
\begin{align}
    \ins{-\laplacian+k^2} \phi &= 4\pi q\delta^3(r)
\end{align}
remembering we are in spherical coordinates so the laplacian operator is 
\begin{equation}
    \laplacian{\phi} = \frac{1}{r^2}\dv{r}\ins{r^2\dv{\phi}{r}}
\end{equation}
finding the derivatives
\begin{align}
    \dv{r}\ins{\frac{q}{r}e^{-kr}} &= q\ins{-\frac{1}{r^2}e^{-kr}-\frac{k}{r}e^{-kr}} = -qe^{-kr}\ins{\frac{1}{r^2}+\frac{k}{r}}\\
    \dv{r}\ins{-qe^{-kr}\ins{1+kr}} &= -q\ins{-ke^{-kr}\ins{1+kr}+ke^{-kr}} = qk^2 re^{-kr}
\end{align}
therefore
\begin{equation}
    \laplacian{\phi} = \frac{qk^2}{r} e^{-kr} = k^2\phi
\end{equation}
and when \(r\neq0\) we get
\begin{equation}
    0 = \ins{k^2-k^2}\phi = 4\pi q\delta^3(r) = 0
\end{equation}
The only point left to check is the origin. Since near the origin \(e^{-kr}\rightarrow 1\), the potential behaves like
\[
\phi \sim \frac{q}{r}.
\]
To see the coefficient of the source, we integrate over a small sphere of radius \(\epsilon\) around the origin:
\[
\int_{V}(-\nabla^2+k^2)\phi\,d^3x.
\]
The \(k^2\phi\) term goes to zero when \(\epsilon\rightarrow0\), because \(\phi\sim q/r\) (so after it's multiplied by the jacobian it goes like \(r\) making the integral of that element tend to 0). Therefore only the Laplacian term contributes:
\[
\int_{V}(-\laplacian\phi)d^3x
=
-\oint_{\partial V}\grad\phi\cdot d\vec S.
\]
Near the origin
\[
\frac{d\phi}{dr}\sim-\frac{q}{r^2},
\]
so
\[
-\oint_{\partial V}\nabla\phi\cdot d\vec S
=
-\left(-\frac{q}{\epsilon^2}\right)4\pi\epsilon^2
=
4\pi q.
\]
So the expression is zero everywhere except the origin, and gives total weight \(4\pi q\) at the origin. Therefore
\[
(-\nabla^2+k^2)\phi=4\pi q\delta^3(r).
\]
Since for a static point charge
\[
J^0=cq\delta^3(r),
\]
we get
\[
\frac{4\pi}{c}J^0=4\pi q\delta^3(r),
\]
so the solution is correct at the origin too.